\secciongrande{INTEGRACIÓN DEL CÓDIGO CON LAS ETIQUETAS HTML: Integración online vs. archivo externo}

\resumen{
  Los navegadores permiten integrar JavaScript directamente en HTML (online) o mediante un archivo externo (\texttt{.js}). La integración externa es la práctica más recomendada por su modularidad, facilidad de depuración y rapidez de carga. Es crucial manejar correctamente las rutas de los archivos y considerar el soporte para navegadores sin JavaScript con la etiqueta \texttt{<noscript>}.
}

\definicion{Integración Online}{Inclusión del código JavaScript directamente dentro de las etiquetas \texttt{<script>} del documento HTML, utilizando el atributo \texttt{type="text/javascript"}.}
\definicion{Integración Externa}{Referencia a un archivo JavaScript separado (con extensión \texttt{.js}) mediante el atributo \texttt{src} dentro de la etiqueta \texttt{<script>}.}
\subtituloA{Métodos de Integración}

\lista{
  **Código Online:** %%VERBATIM\_BLOCK\_0%% |
  **Código Externo:** %%VERBATIM\_BLOCK\_1%%
}
